% Integration-ready prose. Counts and source scope: temporal/results.json.
% This is an extension of the fixed 240-row sample, not a new random sample.
\subsection{Prospective statements, opportunities, and later behavior}
\label{sec:temporal-audit}

We extended the fixed 240-decision sample with a manual reading of every sampled public message and applied private note. The sample consists of three SHA-ranked public-unflagged decisions per player-game, rather than strata defined by whether a reviewer later found the decision noteworthy. Provider summaries and assistant content were excluded from this coding. We retained absent fields, actual legal menus, action provenance, and original source hashes. Fifty-three sampled private fields were absent, all for Opus; none of the 240 actions was an engine default or carried a synthetic fallback note. Consequently, the presence of an explicit plan in this text cannot support a comparison of planning ability across configurations.

The coding found forward-looking own-action passages in 59 records, prospective opponent capacity or outcome passages in 36, and explicit correction or feedback references in 25. These overlapping categories are a screen for temporal claims, not correctness scores. They include contingent goals, possible actions, and unspecified horizons. We assigned every prospective own-action row an endpoint disposition: eight belonged to a complete, text-defined landing-policy subset; nineteen supplied illustrative linked chronologies; thirty-two remained unscorable as single commitments. For the latter, the ledger preserves the reason and subsequent actual opportunities rather than treating noncompletion as failure. One reader performed this semantic coding; the retained source fields and decisions make the classifications inspectable.

\paragraph{A plan requires an opportunity before it can fail.}
Eight sampled records explicitly prioritize buying named bank properties if landed on or on sight. We searched each actor's subsequent actual purchase menus through the end of the game. Six records had no later named-target purchase opportunity. The two remaining cases have different outcomes. In G2, Grok's d161 policy includes Mediterranean Avenue; it buys that property at its later d256 opportunity. In G16, Grok says at d103 to ``buy Virginia on sight, then house pink.'' At d159 it has \$371, Virginia costs \$160, and the delivered menu includes \texttt{buy\_property}. Its applied note still says ``Buy at list is the only way to close pink,'' but the applied action starts an auction. Gemini bids \$380 at d162 and receives Virginia after d164. This is an observed failure to execute the stated purchase at a real opportunity, accompanied by a note/action conflict. Calling it strategic abandonment would misdescribe the record. The six unexercised policies are not failures, and the two repeated G6 policy statements are not independent episodes. We therefore report no adherence percentage.

\paragraph{Execution can preserve a target while changing its route.}
The records contain sustained, realized targets as well as revisions. At G9 d64, Sol targets three houses on each dark blue as funds arrive. It adds a Boardwalk house at d67 and a Park Place house at d102, then reaches three on each at d163, moving from global turn29 to60. This is a realized multi-stage objective. Rolls, transfers and financing intervene, so the chronology does not show that remembering the statement caused the result. A shorter sequence is more tightly specified: G14 d145 states a brown four-house target, and three later building decisions in the same turn reach it, leaving \$619. Eventual target attainment is distinct from preserving the rejected original purchase allocation within a single decision.

The distinction also matters for immediate plans. G20 d99 says to mortgage Vermont next. Instead, d100 first offers it for \$150. After rejection, d102 mortgages it and d103 ends with the anticipated \$157. A next-action test and a same-turn financing-endpoint test would give different descriptions. At G17 d236, a mortgage is explicitly intended to fund three fourth houses; d237 instead ends with a revised reserve argument. At a later changed state, d241 applies one added house, followed by a reserve-based END at d242. This demonstrates revision, while the note's omitted unlocked mortgage capacity prevents treating the revised explanation as evidence of improved financial reasoning.

\paragraph{Opponent assessment is not yet a theory-of-mind score.}
Most candidate opponent forecasts concern public cash, reachable rents, collateral, or hypothetical trade proceeds. These can be tested as financial or movement propositions without inferring what another agent privately believes. At G2 d161, the forecast that withdrawal leaves New York to p2 is realized after two opponents with insufficient cash also drop. This is a correct public-capacity forecast. At G3 d207, the next own management decision still has three bank houses, but follows a doubles continuation before any opponent acts; that first endpoint is not an opponent-behavior test. Stock also remains three at d213 after rival decisions. Neither observation establishes that each opponent could buy no more than the stated number under every legal financing sequence.

At G17 d228, Grok predicts that a rival will outbid its \$99 ceiling and still build. The rival does outbid, jumping to \$100; Grok then faces a \$101 minimum and withdraws. We follow the rival through its d267 elimination. Jail rolls are followed by a fine and debt liquidation before its first management decision at d257 with \$63. No later delivered menu offers building. This is neither an observed deliberate nonbuild nor a falsification of the initial conditional financing claim; alternative financing sequences at those later states were not executed here. In G19 d220, Grok predicts that rivals will not pay it to survive, but no offer is presented before its elimination. The prediction has no observed acceptance test. These cases show why realized outcomes, feasible choices, and attributed motives require different labels. They do not supply calibrated private-belief predictions or a general theory-of-mind measurement.

\paragraph{Revision is observable; learning requires more.}
Within-auction changes offer a useful comparison because ownership and settled auction cost have not yet changed. At G20 d40, Grok gives a numerical ceiling of \$235; d42 tightens it to \$230; d44 bids \$232 and explicitly acknowledges exceeding the newer limit; d46 withdraws. The original \$235 bound is not exceeded. In G17, a roughly \$400 ceiling becomes \$460 and then \$530, with bids of \$410 and \$500, before cash forces withdrawal. The notes offer new rival-resale arguments. These are observed revisions of stated rules, with no demonstrated improvement in value or win probability. Conversely, withdrawing after a rival's \$100 jump in G17 d232 respects the earlier \$99 cash ceiling; it is not evidence of inconsistent bidding.

Together with the separate full-corpus corrective-feedback recurrence audit, these sequences provide bounded answers about adaptation. Agents sometimes revise an explicit strategy, sometimes carry a target through several decisions, and sometimes produce a final action that conflicts with the strategy text at the relevant opportunity. Neither a corrected retry nor a later successful target establishes learned policy. Testing learning would require comparable opportunities, controlled feedback, and outcomes that distinguish improvement from changed assets, rival behavior, or random movement. Testing whether strategic prose reflects effective strategy likewise requires an action-value comparison; this study establishes text--action and opportunity relations rather than global optimality.
